PROMPT PARA IMPLEMENTAR BACKEND A FRONTEND DE PROYECTO YA CREADO:

ANALIZA EL PROYECTO Y REALIZA LO SIGUIENTE:
CREA EL BACKEND CON PHP, EL CUAL ESTARÁ CONECTADO A UNA BASE DE DATOS EN SUPABASE (DIME EN DONDE DEBO LLENAR LOS DATOS PARA LANZAR LA BASE DE DATOS DE DESARROLLO, Y CONECTAR EL CODIGO A ESTA).
ANALIZA EL CODIGO DE FRONTEND PARA CONECTAR ACERTIVAMENTE EL BACKEND, PUEDE TOMAR DE REFERENCIA EL ARCHIVO DE LEEME.md
ORGANIZAR LA ESTRUCTURA DE CARPETAS/ARCHIVOS PARA TENER UN ORDEN VISUAL ENTRE EL FRONT Y EL BACKEND

DEBES CREAR LAS SIGUIENTES TABLAS EN LA BDD (ENTREGAME LAS SENTENCIAS SQLS A EJECUTAR), Y SU RESPECTIVA CLASE Y FUNCIONES (POO) EN EL BACKEND:
USUARIOS
GASTOS_INDISPENSABLES
GASTOS_VARIABLES
PERFIL_FINANCIERO

TODAS LAS TABLAS Y COLUMNAS DEBEN IR EN MAYUSCULA Y CUANDO HAYA ALGUN ESPACIO OCUPAR _

DEBEN EXISTIR HERRAMIENTAS BASICAS PARA OPTIMIZAR EL DESARROLLO, COMO FUNCIONES UTILES COMO DEBUGS, VERIFICACION DE CONEXIÓN CON SUPABASE


DETALLE DE LA TABLA USUARIOS:
    ID_USUARIO
    NOMBRES
    APELLIDOS
    EMAIL_USUARIO
    CONTRASEÑA_USUARIO
    CREATED_AT
    LOGGED_AT

DETALLE PARA TABLA DE PERFIL FINANCIERO:
    ID_PERFILFINANCIERO
    ID_USUARIO
    INGRESO_MENSUAL
    GASTOS_INDISPENSABLES_TOTAL
    AHORRO_BASE
    PERIODO_PRESUPUESTO
    CREATED_AT
    UPDATE_AT;
        
DETALLE PARA LA TABLA GASTOS_INDISPENSABLES:
    ID_GASTOINDISPENSABLE
    ETIQUETA
    MONTO
    CREATED_AT
    
DETALLE PARA LA TABLA GASTOS_VARIABLES:
    ID_GASTOS_VARIABLES
    ETIQUETA
    MONTO
    CATEGORIA
    METODO
    CREATED_AT
